\documentclass[8pt,a4paper]{article}

% ------------------------------------------------
% Packages
% ------------------------------------------------
\usepackage[
    top=0.15in,
    bottom=0.15in,
    left=0.55in,
    right=0.55in
]{geometry}

\usepackage{amsmath,amsfonts,amssymb}
\usepackage{booktabs}
\usepackage{tabularx}
\usepackage{makecell}
\usepackage{caption}
\usepackage{hyperref}
\usepackage{titlesec}

% ------------------------------------------------
% Formatting
% ------------------------------------------------

% Black-and-white hyperlinks
\hypersetup{
    colorlinks=true,
    allcolors=black
}

% Reduce line spacing
\linespread{0.52}

% Reduce paragraph spacing
\setlength{\parindent}{0pt}
\setlength{\parskip}{1pt}

% Reduce section heading spacing
\titlespacing*{\subsection}
{0pt}{7pt}{3pt}

% Table formatting
\setlength{\tabcolsep}{5pt}
\renewcommand{\arraystretch}{1.05}

\captionsetup[table]{
    font=small,
    labelfont=bf,
    position=top,
    skip=3pt
}

% ------------------------------------------------
% Document
% ------------------------------------------------
\begin{document}
\thispagestyle{empty}

% ------------------------------------------------
% Title and Team Information
% ------------------------------------------------
\begin{center}

{\large\bfseries
Similarity-Cue Injection for Flexible Frame Selection in Video Question Answering}\\[4pt]

\textbf{Team:} Mitra Mandal
\quad | \quad
\textbf{Members:}
Abhishek Kumar Jha (230106001), Aditya Dhaniyaal (230106002),\\[2pt]
Dev Shukla (230106020), Dhruv Mishra (230106022), Manish Hingar (230106040)\\[3pt]

\textbf{GitHub:}
\url{https://github.com/manishhingar28/DLCV_COURSE_PROJECT}

\end{center}

\vspace{-5pt}

% ------------------------------------------------
\subsection*{I. Abstract (Maximum 100 Words)}
% ------------------------------------------------
\begin{enumerate}
    \item Video-language models are costly because every input frame passes through a large model.
    \item So systems feed a few uniform frames, which can miss the key moment or waste budget.
    \item Buch et al.\ (CVPR 2025) place a light Flexible Frame Selector (FFS) before the large model: it scores frames against the question, keeps the useful ones via a differentiable top-$k$, and uses learnable \texttt{[CTRL]} tokens to decide \emph{how many} ($k{-}c$ to $k$).
    \item We read this as learning, from the answer loss alone, which evidence a question needs, so fewer frames match dense sampling.
\end{enumerate}

% ------------------------------------------------
\subsection*{II. Proposed Novelty \& Technical Contribution}
% ------------------------------------------------
\begin{enumerate}
    \item \textbf{Similarity-cue injection} (new module, 0.07M params): FFS learns frame relevance from the answer loss alone through a straight-through top-$k$, a weak signal with frozen EVA02-CLIP-L/14 features (its top-4 is only +0.7 over uniform). We compute four \emph{zero-shot relevance cues} per frame from the same features (cosine to the question, best and second-best ``question\,+\,option'' similarity, their gap), standardise them per video, and feed them to the frame tokens and, via a linear skip, to the selection logit.
    \item Type: architectural modification (new cue module at the selector input); encoder, \texttt{[CTRL]}, Gumbel top-$k$, losses and downstream model unchanged; no labels, no extra network.
    \item Difference: FFS learns question--frame similarity through attention from scratch; ours starts from a zero-shot ranker and learns a correction. Both share our FFS recipe (\texttt{[CTRL]} count hinge from fixed-$k$ policies).
\end{enumerate}

% ------------------------------------------------
\subsection*{III. Theoretical / Mathematical Justification}
% ------------------------------------------------
\begin{enumerate}
    \item $s_{mt}=\langle\hat t_m,\hat f_t\rangle,\quad c_t=\big[\langle\hat q,\hat f_t\rangle,\ s_{(1)t},\ s_{(2)t},\ s_{(1)t}-s_{(2)t}\big],\quad \tilde c_t=(c_t-\mu_{\mathrm{vid}})/\sigma_{\mathrm{vid}},\quad l_t=g_\theta\big(h_t+\phi(\tilde c_t)\big)+w^\top\tilde c_t.$
    \item $\hat f_t,\hat q,\hat t_m$: normalised frame, question and ``question\,+\,option $m$'' features; $s_{(1)},s_{(2)}$: two best options; $h_t$: frame token; $g_\theta$: FFS encoder/head; $l_t$: selection logit. $w{=}0,\phi{=}0$ gives FFS, so our class contains the baseline. \textbf{Residual learning:} without cues, relevance must be learned as a bilinear map $\langle W_Q\hat q,W_K\hat f_t\rangle$ in 768-d from noisy straight-through gradients; the cues give it in closed form ($g_\theta{=}0$ ranks by $w^\top\tilde c_t$), so $g_\theta$ learns only a correction (lower sample complexity). \textbf{Margin = evidence:} with $p(m|f_t)\propto e^{s_{mt}/T}$, a large $s_{(1)t}-s_{(2)t}$ means low entropy $H(a|q,f_t)$: a frame that \emph{discriminates} answers, found without the label. \textbf{Standardisation:} top-$k$ depends only on the order of $l_t$, unchanged by per-video affine shifts, while CLIP cosines shift strongly across videos (mean 0.89); $\tilde c_t$ removes this nuisance.
    \item Causal link: better ranking $\Rightarrow$ the same $k$ frames carry more answer evidence $\Rightarrow$ higher accuracy at equal cost (tested in Table~\ref{tab:results}). Ablation (RAS): greedy $i_t=\arg\max_i l_i-\beta\sum_{j\in A_{t-1}}S_{ij}$ (submodular), assuming near-duplicate frames waste budget.
\end{enumerate}

% ------------------------------------------------
\subsection*{IV. Comparative Results \& Empirical Evaluation}
% ------------------------------------------------
\begin{enumerate}
    \item NExT-QA multiple choice, val (4,996 Qs; train 34,132), 32-frame pool, frozen EVA02-CLIP-L/14 features at 224\,px; a small frozen QA model (2.0M params) replaces the paper's 8B VLM, so absolute accuracy is not comparable. Shared downstream model and recipe; tuning on held-out \emph{training} videos only; 3 seeds.
    \item Reproduced baseline (FFS-style input) vs ours at 4 frames and with adaptive FFS (1--4); Frames = frames given to the downstream model.
    \item Takeaway: +1.05 accuracy at 4 frames and +0.85 at 1 frame for +0.07M parameters; with cues, 4 selected frames beat all 32 uniform frames at 1/8 of the downstream cost.
\end{enumerate}

% ------------------------------------------------
% Results Table
% ------------------------------------------------
\begin{center}

\captionof{table}{
Quantitative comparison of the baseline and
proposed method on NExT-QA (val, 3 seeds).
Best results are highlighted in bold.
}

\label{tab:results}

\small

\begin{tabularx}{\linewidth}{
    @{}X c c c c c c@{}
}

\toprule

\textbf{Method}
&
\makecell{\textbf{Acc}\\(\% $\uparrow$)}
&
\makecell{\textbf{Causal}\\($\uparrow$)}
&
\makecell{\textbf{Temporal}\\($\uparrow$)}
&
\makecell{\textbf{Descr.}\\($\uparrow$)}
&
\makecell{\textbf{Frames}\\($\downarrow$)}
&
\makecell{\textbf{Params}\\(M)}
\\

\midrule

Uniform sampling, 32 frames / 4 frames
& 48.40 / 47.94 & 46.9 / 46.8 & 46.3 / 45.1 & 57.9 / 57.7 & 32 / 4 & -- \\

Baseline: Top-4 (FFS-style input)
& 48.66 $\pm$ 0.19 & 47.8 & 45.7 & 57.7 & 4 & 2.05 \\

\textbf{Ours: Top-4 + similarity cues}
& \textbf{49.71 $\pm$ 0.10} & \textbf{48.3} & \textbf{47.4} & \textbf{59.1} & 4 & 2.12 \\

Baseline: FFS [paper, re-impl.], adaptive 1--4
& 46.35 $\pm$ 0.19 & 44.9 & 44.0 & 56.2 & \textbf{1.0} & 2.05 \\

\textbf{Ours: FFS + similarity cues}, adaptive
& 47.20 $\pm$ 0.10 & 45.4 & 45.0 & 57.6 & \textbf{1.0} & 2.12 \\

Ablation: Top-4 + cues + RAS
& 49.71 $\pm$ 0.10 & 48.3 & 47.3 & 59.3 & 4 & 2.12 \\

\midrule

$\Delta$ (Ours -- Baseline), top-4 / FFS
& +1.05 / +0.85 & +0.5 / +0.6 & +1.7 / +1.0 & +1.4 / +1.4 & same & +0.07 \\

\bottomrule

\end{tabularx}

\end{center}

\vspace{-5pt}

\textbf{Result Interpretation:}
Cues improve both selectors in every question type (most: Temporal) by about 5$\times$ the seed std at the same cost; with cues, 4 selected frames \emph{beat} all 32 uniform frames (+1.3) at 1/8 of the cost, while the FFS-style input only matches them (+0.3), confirming that cues improve \emph{which} frames are chosen. Trade-offs: adaptive selectors settled on one frame (this downstream model gains only 3.0 points from 1$\to$32 frames); RAS left picks unchanged, and learned picks are \emph{more} similar than uniform ones (0.92 vs 0.89) yet more accurate, so evidence clusters around key moments. Caveats: no per-question significance test; small downstream model; selector reads all 32 frames' features.

\end{document}